% ============================================================
%  Background chapter: Kinematics of the ViperX-300
%  Draft — editable working file
%  Include in main thesis with % ============================================================
%  Background chapter: Kinematics of the ViperX-300
%  Draft — editable working file
%  Include in main thesis with % ============================================================
%  Background chapter: Kinematics of the ViperX-300
%  Draft — editable working file
%  Include in main thesis with % ============================================================
%  Background chapter: Kinematics of the ViperX-300
%  Draft — editable working file
%  Include in main thesis with \input{kinematics_background}
% ============================================================

Robotics may come in many different forms but they all have something in
common: they share the same basics. The modern studies of robots include
everything from kinematics and trajectory generation to motion planning. In
order to grasp these more advanced concepts one must first recognise the
fundamentals. As a foundation of the subject, the following draws on
\cite{lynch2017modern}, which provides a clear explanation of the key concepts
of robotic manipulators.

\subsection{Configuration Space}
%When entering the mathematical way of explaining how manipulators move and  (Explaining joint space vs. task space)
To construct a robot you need various different components. The basic \textit{link} is rigid body, to permit movement of one link relative to the next one, a \textit{joint} is connected between the two consecutive links. At its core, a manipulator is a chain of such links coupled by joints. To realise the complex current pose of the manipulator, it can always be reduced to a couple of number describing the lengths of the links and the angle between them.
The complete specification of every position of each point of the robot is a \textit{configuration}. The minimum number of real-valued coordinates required to specify the configuration is the number of degrees of freedom of the robot, and the space of all attainable configurations is the \textit{configuration space}, of C-space \cite{lynch2017modern}. For a manipulator built from any number of $n$ joints adding one degree of freedom each, a configuration is conveniently collected intro the joint variable vector
\begin{equation}
    \mathbf{q} = \begin{bmatrix} q_1 & q_2 & \cdots & q_n \end{bmatrix}^\top \in \mathbb{R}^n,
\end{equation}
where each $q_i$ is a joint angle for a revolute joint or a linear displacement for prismatic joint.
When addressing the configuration space, generally the space is not configured to the conventional Euclidean space but rather a manifold. Taking a single revolute joint for example, it cant only move along a circle's outer rim and so forth its C-space is confined to a circle. The configuration space of a $n$-joint revolute arm will therefore result in an $n$-torus.

The configuration space described above, with each point given by the joint
variable vector $\mathbf{q}$, is what is referred to as the \textit{joint space}
of the manipulator. These are the variables that are directly driven by the
actuators and read back by the sensors. The joint space is however seldom the
space in which a task is actually described. When a manipulator is asked to reach
a point or to align its gripper, the task is expressed through the position and
orientation of the end-effector, which belong to a separate space called the
\textit{task space} (or operational space). Rather than joint values, the task
space is built from poses in ordinary three-dimensional space.

The link between the two spaces is the kinematics of the manipulator. The
\textit{forward kinematics} takes a set of joint values from the joint space and
returns the corresponding end-effector pose in the task space, while the
\textit{inverse kinematics} goes the other way, starting from a desired pose and
solving for the joint values that achieve it. Deriving and studying these two
mappings is the main focus of the sections that follow.

\subsection{Degrees of Freedom}

When specifying the \gls{dof} of a robot, one is specifying how the robot can
move in space. To start, the classic but straightforward example given in
\cite{lynch2017modern} explains the simple door fixed by hinges to the frame.
On its own, a rigid body moving freely in three-dimensional space has six
degrees of freedom. Attaching the door to the wall through the hinge introduces
constraints: the hinge permits rotation about its axis through an angle $\theta$
while removing every other possible motion. Five of the six degrees of freedom
are thereby constrained, leaving a single independent coordinate. The door
therefore has only one degree of freedom.

\begin{figure}[ht]
  \centering
  \input{figures/joints_dof.tex}
  \caption{The 6 different joint categories.}
  \label{fig:joints_dof}
\end{figure}

\subsection{Anatomy of an Open-Chain Robot}

The manipulators considered in this thesis are \textit{open-chain}, or serial,
mechanisms. In such a structure the links form a single unbranched sequence:
starting from a fixed \textit{base} that is rigidly attached to the environment,
each link is connected to the next by exactly one joint, terminating in the
\textit{end-effector}, the tool or gripper through which the robot interacts with
its surroundings. The defining property of an open chain is that there is one and
only one kinematic path from the base to the end-effector, in contrast to
closed-chain mechanisms where links form one or more loops. This single-path
structure is what makes the recursive construction of the kinematics in the
following sections possible.

The \textit{joints} are the elements that introduce motion into the chain, and
each one removes some of the relative freedom between the two links it connects
while permitting motion along the remaining directions. The two joint types
relevant here are the \textit{revolute} joint, which allows a pure rotation about
a fixed axis and is parameterised by an angle, and the \textit{prismatic} joint,
which allows a pure translation along a fixed axis and is parameterised by a
displacement. Both contribute a single degree of freedom. More general joints
such as cylindrical, spherical and universal joints exist and can be modelled as
combinations of these elementary joints; the six joint categories are illustrated
in Figure~\ref{fig:joints_dof}. The ViperX-300 studied in this work is a 6-DOF
open-chain manipulator constructed entirely from revolute joints, a configuration
that grants its end-effector six degrees of freedom and thus, in principle, the
ability to attain an arbitrary position \textit{and} orientation within its
reachable workspace.

Motion is produced by the \textit{actuators} that drive the joints. In modern
manipulators these are most commonly electric servomotors, typically brushless DC
or stepper motors combined with a gear reduction to trade speed for the torque
required to support the arm and its payload; hydraulic and pneumatic actuators are
also used where high force density or compliance is needed. Each actuator applies
a torque (for a revolute joint) or a force (for a prismatic joint) at its
respective joint in response to a commanded control signal. Because the links are
coupled, a torque applied at one joint propagates through the entire chain and
influences the motion of every link beyond it. Capturing this coupling
quantitatively is precisely the role of the dynamic model, which relates the
vector of joint torques to the resulting joint accelerations and is treated later
in this chapter.

\subsection{Rotation Matrix}

One of the fundamental tasks in modelling an open-chain robotic manipulator is to
describe the position and orientation of its end-effector. To do this, a
\textit{coordinate frame} is attached to each rigid body of the robot, giving a
systematic way of relating points and directions along the chain. The kinematics
then describes the geometric relations between these frames. To relate one frame
to another, two separate pieces of information are needed: the \textit{position}
of its origin, described by a translation vector, and its \textit{orientation},
described by a rotation matrix.

The orientation of a frame $B$ relative to a frame $A$ is conveniently described
by attaching an orthonormal frame to the body and expressing its unit vectors in
the coordinates of the reference frame \cite{siciliano2009robotics}. Collecting
the three unit vectors $\mathbf{x}_B$, $\mathbf{y}_B$ and $\mathbf{z}_B$ as
columns gives the $3 \times 3$ \textit{rotation matrix}
\begin{equation}
    ^A R_B = \begin{bmatrix}
        \mathbf{x}_B & \mathbf{y}_B & \mathbf{z}_B
    \end{bmatrix},
    \label{eq:rotation_matrix}
\end{equation}
whose entries are the direction cosines of the axes of frame $B$ with respect to
the axes of frame $A$.

Since the columns of $^A R_B$ are the unit vectors of an orthonormal frame, they
are mutually orthogonal and of unit norm. These conditions are compactly expressed
as
\begin{equation}
    ^A R_B^\top \, ^A R_B = I,
    \label{eq:orthogonality}
\end{equation}
that is, the rotation matrix is orthogonal. Postmultiplying \eqref{eq:orthogonality}
by the inverse matrix gives the useful result
\begin{equation}
    ^A R_B^{-1} = {}^A R_B^\top = {}^B R_A,
\end{equation}
so that reversing the roles of the two frames requires nothing more than a
transposition. The determinant of a rotation matrix equals $+1$ for a right-handed
frame and $-1$ for a left-handed one; since only right-handed frames are used here,
the matrices of interest belong to the \textit{special orthogonal group}
\begin{equation}
    SO(3) = \left\{ R \in \mathbb{R}^{3\times 3} \;\middle|\;
    R^\top R = I, \; \det R = +1 \right\}.
\end{equation}

Frames obtained by rotating the reference frame about one of its coordinate axes
are described by the \textit{elementary rotation matrices}. Taking the rotation
positive when made counter-clockwise about the axis in question, a rotation by an
angle $\vartheta$ about the $z$-, $y$- and $x$-axis is respectively given by
\begin{equation}
    R_z(\vartheta) = \begin{bmatrix}
        \cos\vartheta & -\sin\vartheta & 0 \\
        \sin\vartheta & \cos\vartheta & 0 \\
        0 & 0 & 1
    \end{bmatrix}, \quad
    R_y(\vartheta) = \begin{bmatrix}
        \cos\vartheta & 0 & \sin\vartheta \\
        0 & 1 & 0 \\
        -\sin\vartheta & 0 & \cos\vartheta
    \end{bmatrix}, \quad
    R_x(\vartheta) = \begin{bmatrix}
        1 & 0 & 0 \\
        0 & \cos\vartheta & -\sin\vartheta \\
        0 & \sin\vartheta & \cos\vartheta
    \end{bmatrix}.
    \label{eq:elementary_rotations}
\end{equation}
For these matrices $R_k(-\vartheta) = R_k^\top(\vartheta)$ holds for
$k = x, y, z$. They serve as the building blocks from which rotations about an
arbitrary axis in space are composed, and they are the ones that reappear in the
Denavit-Hartenberg convention.

Beyond describing orientation, the rotation matrix also relates the coordinates
of a point in two frames sharing a common origin. If a point $P$ is expressed as
$^B\mathbf{p}$ in frame $B$, its coordinates in frame $A$ follow from
\eqref{eq:rotation_matrix} as
\begin{equation}
    ^A\mathbf{p} = {}^A R_B \, ^B\mathbf{p},
    \label{eq:coordinate_transform}
\end{equation}
with the inverse transformation given by
$^B\mathbf{p} = {}^A R_B^\top \, ^A\mathbf{p}$. A third interpretation is obtained
by regarding $^A R_B$ as an operator acting within a single frame: the product
$R\mathbf{p}$ yields a vector rotated with respect to $\mathbf{p}$ but of the same
norm, as follows directly from the orthogonality property
\eqref{eq:orthogonality}. A rotation matrix therefore admits three equivalent
geometrical meanings \cite{siciliano2009robotics}: it describes the mutual
orientation of two coordinate frames, it transforms the coordinates of a point
between two frames with a common origin, and it rotates a vector within one frame.

Rotations are combined by matrix multiplication. Considering three frames $0$, $1$
and $2$ with a common origin, applying \eqref{eq:coordinate_transform} twice gives
\begin{equation}
    ^0 R_2 = {}^0 R_1 \, ^1 R_2,
    \label{eq:rotation_composition}
\end{equation}
which is interpreted as a sequence of rotations, each defined with respect to the
frame resulting from the preceding one. The frame about whose axes a rotation is
made is termed the \textit{current frame}, and composition about current frames is
accordingly obtained by \textit{postmultiplication} in the order the rotations are
performed. If instead every rotation is referred to the initial frame, which is
then termed the \textit{fixed frame}, the same reasoning leads to composition by
\textit{premultiplication}. Since matrix multiplication is not commutative, two
rotations in general do not commute, and the resulting orientation depends on the
order in which they are applied. The Denavit-Hartenberg convention introduced
below is built on composition about current frames, and hence on postmultiplication
of the elementary rotations in \eqref{eq:elementary_rotations}.

\subsection{Homogeneous Transformation Matrix}

The rotation matrix and the translation vector can be combined into a single
description of \textit{pose} by means of the \textit{homogeneous transformation
matrix}, which allows the pose of an entire kinematic chain to be computed by
ordinary matrix multiplication. The pose of frame $B$
relative to frame $A$ is written as the $4 \times 4$ matrix
\begin{equation}
^A T_B = \begin{bmatrix}
n_x & s_x & a_x & P_x \\
n_y & s_y & a_y & P_y \\
n_z & s_z & a_z & P_z \\
0 & 0 & 0 & 1
\end{bmatrix} = \begin{bmatrix}
^A R_B & ^A P_B \\
\mathbf{0}^\top & 1
\end{bmatrix},
\end{equation}
where the upper-left $3 \times 3$ block $^A R_B$ is the rotation matrix describing
the relative orientation and the upper-right column $^A P_B$ is the position of the
origin of frame $B$ expressed in frame $A$. The columns of the rotation block are
frequently denoted $\mathbf{n}$, $\mathbf{s}$ and $\mathbf{a}$ (normal, sliding and
approach), a notation that originates from the description of gripper orientation.
The bottom row $\begin{bmatrix} 0 & 0 & 0 & 1 \end{bmatrix}$ has no geometric content
of its own but is what makes the composition of transformations reduce to matrix
multiplication.

The set of all such matrices forms the \textit{special Euclidean group} $SE(3)$, the
group of rigid-body motions in three-dimensional space. Its two central properties
are exactly the ones that make it the natural tool for serial manipulators. First,
transformations \textit{compose}: if the pose of frame $C$ is known relative to $B$,
and that of $B$ relative to $A$, then the pose of $C$ relative to $A$ is obtained
simply as
\begin{equation}
    ^A T_C = {}^A T_B \, ^B T_C,
\end{equation}
which generalises directly to a chain of arbitrarily many frames. Second, the inverse
transformation has the closed form
\begin{equation}
    ^A T_B^{-1} = {}^B T_A = \begin{bmatrix}
        ^A R_B^\top & -\,^A R_B^\top \, ^A P_B \\
        \mathbf{0}^\top & 1
    \end{bmatrix},
\end{equation}
so that changing the direction of a transformation requires no numerical matrix
inversion. The compositional property is the cornerstone of forward kinematics: by
attaching a frame to every link and writing the transformation between consecutive
frames, the pose of the end-effector relative to the base is recovered by multiplying
the individual link transformations together. The systematic assignment of these link
frames is governed by the Denavit-Hartenberg convention described next.

\subsection{The Denavit-Hartenberg (D-H) Convention}

The \acrshort{dh} convention is a standardised procedure for attaching coordinate
frames to the links of a serial manipulator. Without such a standard, the placement
of link frames would be arbitrary and the transformation between two consecutive
frames would in general require six parameters --- three for position and three for
orientation. The insight underlying the D-H convention is that, by choosing the frames
according to a fixed set of rules, the relative pose of two consecutive frames can
always be described by only \textit{four} parameters instead of six
\cite{siciliano2009robotics}. This reduction is achieved by aligning each frame's
$z$-axis with its joint axis and its $x$-axis along the common normal between
consecutive joint axes, which removes two of the degrees of freedom in the frame
placement.

The four parameters, illustrated for link $i$, are the \textit{link length} $a_i$,
the \textit{link twist} $\alpha_i$, the \textit{link offset} $d_i$, and the
\textit{joint angle} $\vartheta_i$. For a revolute joint the joint angle
$\vartheta_i$ is the variable quantity and the remaining three parameters are fixed
by the geometry of the link; for a prismatic joint the offset $d_i$ is variable and
the other three are constant.

The convention works along a recursive algorithm that transforms frame $i{-}1$ into
frame $i$ through four successive elementary transformations, performed in order:
\begin{enumerate}
    \item a rotation about the $z_{i-1}$-axis by the angle $\theta_i$, aligning
          $x_{i-1}$ with $x_i$;
    \item a translation along the $z_{i-1}$-axis by the offset $d_i$;
    \item a translation along the (now rotated) $x_i$-axis by the length $a_i$;
    \item a rotation about the $x_i$-axis by the twist $\alpha_i$, aligning
          $z_{i-1}$ with $z_i$.
\end{enumerate}
Expressing each step as a homogeneous transformation and multiplying them together
yields the link transformation matrix
\begin{equation}
^{i-1}T_i = \begin{bmatrix}
\cos\theta_i & -\sin\theta_i\cos\alpha_i & \sin\theta_i\sin\alpha_i & a_i\cos\theta_i \\
\sin\theta_i & \cos\theta_i\cos\alpha_i & -\cos\theta_i\sin\alpha_i & a_i\sin\theta_i \\
0 & \sin\alpha_i & \cos\alpha_i & d_i \\
0 & 0 & 0 & 1
\end{bmatrix}.
\label{eq:dh_matrix}
\end{equation}
Applying this procedure to every joint of the ViperX-300 yields the parameter set
listed in Table~\ref{tab:dh_params}, from which the complete kinematic model of the
manipulator follows.

\subsection{Forward and Inverse Kinematics}

\subsubsection{Forward Kinematics}
With a link frame assigned to each joint and the corresponding D-H transformation
\eqref{eq:dh_matrix} written for every link, the \textit{forward kinematics} of the
manipulator follows directly from the compositional property of homogeneous
transformations. Given a joint configuration
$\mathbf{q} = \begin{bmatrix} q_1 & \cdots & q_n \end{bmatrix}^\top$, the pose of the
end-effector frame $n$ relative to the base frame $0$ is obtained by chaining the
individual link transformations:
\begin{equation}
    ^0 T_n(\mathbf{q}) = {}^0 T_1(q_1)\, ^1 T_2(q_2) \;\cdots\; ^{n-1} T_n(q_n).
    \label{eq:fk_chain}
\end{equation}
The result is a single $4 \times 4$ matrix encoding both the position and the
orientation of the end-effector as a function of the joint variables. Forward
kinematics is, from a computational standpoint, entirely benign: for any given
$\mathbf{q}$ the pose is unique and is obtained by a fixed sequence of matrix
multiplications, with no iteration or ambiguity. It constitutes the mapping from
joint space to task space introduced earlier.

\subsubsection{Inverse Kinematics}
The \textit{inverse kinematics} problem reverses this mapping: given a desired
end-effector pose $^0 T_n^{\,\star}$, the task is to find a joint configuration
$\mathbf{q}$ for which $^0 T_n(\mathbf{q}) = {}^0 T_n^{\,\star}$. Although it is the
inverse of a problem that was straightforward, it is considerably more delicate, for
three reasons. First, \textit{existence} is not guaranteed: a solution exists only if
the requested pose lies within the reachable workspace of the manipulator. Second, the
solution is generally \textit{not unique}: a 6-DOF arm can frequently reach the same
pose through several distinct configurations --- the familiar elbow-up and elbow-down
postures being the classic example --- so the inverse map is one-to-many. Third, the
equations to be solved are \textit{nonlinear} in the joint variables, since the entries
of \eqref{eq:fk_chain} contain products of sines and cosines of the joint angles, and
for a general manipulator no closed-form solution exists.

In such cases the inverse kinematics must be solved numerically, typically by iterative
schemes based on the manipulator Jacobian, which raise the usual concerns of
convergence, computational cost and behaviour near singularities. Closed-form analytic
solutions are nonetheless available for manipulators whose geometry satisfies certain
conditions; the most important of these is the presence of a \textit{spherical wrist},
in which the axes of the final three joints intersect at a common point. This geometry
allows the problem to be \textit{decoupled}, with the wrist centre position determining
the first three joints and the end-effector orientation determining the last three,
each subproblem then admitting an analytic solution. Whether such a decoupling applies
to the ViperX-300 follows directly from its D-H parameters and is examined in the
kinematic analysis of this work.

% ============================================================

Robotics may come in many different forms but they all have something in
common: they share the same basics. The modern studies of robots include
everything from kinematics and trajectory generation to motion planning. In
order to grasp these more advanced concepts one must first recognise the
fundamentals. As a foundation of the subject, the following draws on
\cite{lynch2017modern}, which provides a clear explanation of the key concepts
of robotic manipulators.

\subsection{Configuration Space}
%When entering the mathematical way of explaining how manipulators move and  (Explaining joint space vs. task space)
To construct a robot you need various different components. The basic \textit{link} is rigid body, to permit movement of one link relative to the next one, a \textit{joint} is connected between the two consecutive links. At its core, a manipulator is a chain of such links coupled by joints. To realise the complex current pose of the manipulator, it can always be reduced to a couple of number describing the lengths of the links and the angle between them.
The complete specification of every position of each point of the robot is a \textit{configuration}. The minimum number of real-valued coordinates required to specify the configuration is the number of degrees of freedom of the robot, and the space of all attainable configurations is the \textit{configuration space}, of C-space \cite{lynch2017modern}. For a manipulator built from any number of $n$ joints adding one degree of freedom each, a configuration is conveniently collected intro the joint variable vector
\begin{equation}
    \mathbf{q} = \begin{bmatrix} q_1 & q_2 & \cdots & q_n \end{bmatrix}^\top \in \mathbb{R}^n,
\end{equation}
where each $q_i$ is a joint angle for a revolute joint or a linear displacement for prismatic joint.
When addressing the configuration space, generally the space is not configured to the conventional Euclidean space but rather a manifold. Taking a single revolute joint for example, it cant only move along a circle's outer rim and so forth its C-space is confined to a circle. The configuration space of a $n$-joint revolute arm will therefore result in an $n$-torus.

The configuration space described above, with each point given by the joint
variable vector $\mathbf{q}$, is what is referred to as the \textit{joint space}
of the manipulator. These are the variables that are directly driven by the
actuators and read back by the sensors. The joint space is however seldom the
space in which a task is actually described. When a manipulator is asked to reach
a point or to align its gripper, the task is expressed through the position and
orientation of the end-effector, which belong to a separate space called the
\textit{task space} (or operational space). Rather than joint values, the task
space is built from poses in ordinary three-dimensional space.

The link between the two spaces is the kinematics of the manipulator. The
\textit{forward kinematics} takes a set of joint values from the joint space and
returns the corresponding end-effector pose in the task space, while the
\textit{inverse kinematics} goes the other way, starting from a desired pose and
solving for the joint values that achieve it. Deriving and studying these two
mappings is the main focus of the sections that follow.

\subsection{Degrees of Freedom}

When specifying the \gls{dof} of a robot, one is specifying how the robot can
move in space. To start, the classic but straightforward example given in
\cite{lynch2017modern} explains the simple door fixed by hinges to the frame.
On its own, a rigid body moving freely in three-dimensional space has six
degrees of freedom. Attaching the door to the wall through the hinge introduces
constraints: the hinge permits rotation about its axis through an angle $\theta$
while removing every other possible motion. Five of the six degrees of freedom
are thereby constrained, leaving a single independent coordinate. The door
therefore has only one degree of freedom.

\begin{figure}[ht]
  \centering
  \input{figures/joints_dof.tex}
  \caption{The 6 different joint categories.}
  \label{fig:joints_dof}
\end{figure}

\subsection{Anatomy of an Open-Chain Robot}

The manipulators considered in this thesis are \textit{open-chain}, or serial,
mechanisms. In such a structure the links form a single unbranched sequence:
starting from a fixed \textit{base} that is rigidly attached to the environment,
each link is connected to the next by exactly one joint, terminating in the
\textit{end-effector}, the tool or gripper through which the robot interacts with
its surroundings. The defining property of an open chain is that there is one and
only one kinematic path from the base to the end-effector, in contrast to
closed-chain mechanisms where links form one or more loops. This single-path
structure is what makes the recursive construction of the kinematics in the
following sections possible.

The \textit{joints} are the elements that introduce motion into the chain, and
each one removes some of the relative freedom between the two links it connects
while permitting motion along the remaining directions. The two joint types
relevant here are the \textit{revolute} joint, which allows a pure rotation about
a fixed axis and is parameterised by an angle, and the \textit{prismatic} joint,
which allows a pure translation along a fixed axis and is parameterised by a
displacement. Both contribute a single degree of freedom. More general joints
such as cylindrical, spherical and universal joints exist and can be modelled as
combinations of these elementary joints; the six joint categories are illustrated
in Figure~\ref{fig:joints_dof}. The ViperX-300 studied in this work is a 6-DOF
open-chain manipulator constructed entirely from revolute joints, a configuration
that grants its end-effector six degrees of freedom and thus, in principle, the
ability to attain an arbitrary position \textit{and} orientation within its
reachable workspace.

Motion is produced by the \textit{actuators} that drive the joints. In modern
manipulators these are most commonly electric servomotors, typically brushless DC
or stepper motors combined with a gear reduction to trade speed for the torque
required to support the arm and its payload; hydraulic and pneumatic actuators are
also used where high force density or compliance is needed. Each actuator applies
a torque (for a revolute joint) or a force (for a prismatic joint) at its
respective joint in response to a commanded control signal. Because the links are
coupled, a torque applied at one joint propagates through the entire chain and
influences the motion of every link beyond it. Capturing this coupling
quantitatively is precisely the role of the dynamic model, which relates the
vector of joint torques to the resulting joint accelerations and is treated later
in this chapter.

\subsection{Rotation Matrix}

One of the fundamental tasks in modelling an open-chain robotic manipulator is to
describe the position and orientation of its end-effector. To do this, a
\textit{coordinate frame} is attached to each rigid body of the robot, giving a
systematic way of relating points and directions along the chain. The kinematics
then describes the geometric relations between these frames. To relate one frame
to another, two separate pieces of information are needed: the \textit{position}
of its origin, described by a translation vector, and its \textit{orientation},
described by a rotation matrix.

The orientation of a frame $B$ relative to a frame $A$ is conveniently described
by attaching an orthonormal frame to the body and expressing its unit vectors in
the coordinates of the reference frame \cite{siciliano2009robotics}. Collecting
the three unit vectors $\mathbf{x}_B$, $\mathbf{y}_B$ and $\mathbf{z}_B$ as
columns gives the $3 \times 3$ \textit{rotation matrix}
\begin{equation}
    ^A R_B = \begin{bmatrix}
        \mathbf{x}_B & \mathbf{y}_B & \mathbf{z}_B
    \end{bmatrix},
    \label{eq:rotation_matrix}
\end{equation}
whose entries are the direction cosines of the axes of frame $B$ with respect to
the axes of frame $A$.

Since the columns of $^A R_B$ are the unit vectors of an orthonormal frame, they
are mutually orthogonal and of unit norm. These conditions are compactly expressed
as
\begin{equation}
    ^A R_B^\top \, ^A R_B = I,
    \label{eq:orthogonality}
\end{equation}
that is, the rotation matrix is orthogonal. Postmultiplying \eqref{eq:orthogonality}
by the inverse matrix gives the useful result
\begin{equation}
    ^A R_B^{-1} = {}^A R_B^\top = {}^B R_A,
\end{equation}
so that reversing the roles of the two frames requires nothing more than a
transposition. The determinant of a rotation matrix equals $+1$ for a right-handed
frame and $-1$ for a left-handed one; since only right-handed frames are used here,
the matrices of interest belong to the \textit{special orthogonal group}
\begin{equation}
    SO(3) = \left\{ R \in \mathbb{R}^{3\times 3} \;\middle|\;
    R^\top R = I, \; \det R = +1 \right\}.
\end{equation}

Frames obtained by rotating the reference frame about one of its coordinate axes
are described by the \textit{elementary rotation matrices}. Taking the rotation
positive when made counter-clockwise about the axis in question, a rotation by an
angle $\vartheta$ about the $z$-, $y$- and $x$-axis is respectively given by
\begin{equation}
    R_z(\vartheta) = \begin{bmatrix}
        \cos\vartheta & -\sin\vartheta & 0 \\
        \sin\vartheta & \cos\vartheta & 0 \\
        0 & 0 & 1
    \end{bmatrix}, \quad
    R_y(\vartheta) = \begin{bmatrix}
        \cos\vartheta & 0 & \sin\vartheta \\
        0 & 1 & 0 \\
        -\sin\vartheta & 0 & \cos\vartheta
    \end{bmatrix}, \quad
    R_x(\vartheta) = \begin{bmatrix}
        1 & 0 & 0 \\
        0 & \cos\vartheta & -\sin\vartheta \\
        0 & \sin\vartheta & \cos\vartheta
    \end{bmatrix}.
    \label{eq:elementary_rotations}
\end{equation}
For these matrices $R_k(-\vartheta) = R_k^\top(\vartheta)$ holds for
$k = x, y, z$. They serve as the building blocks from which rotations about an
arbitrary axis in space are composed, and they are the ones that reappear in the
Denavit-Hartenberg convention.

Beyond describing orientation, the rotation matrix also relates the coordinates
of a point in two frames sharing a common origin. If a point $P$ is expressed as
$^B\mathbf{p}$ in frame $B$, its coordinates in frame $A$ follow from
\eqref{eq:rotation_matrix} as
\begin{equation}
    ^A\mathbf{p} = {}^A R_B \, ^B\mathbf{p},
    \label{eq:coordinate_transform}
\end{equation}
with the inverse transformation given by
$^B\mathbf{p} = {}^A R_B^\top \, ^A\mathbf{p}$. A third interpretation is obtained
by regarding $^A R_B$ as an operator acting within a single frame: the product
$R\mathbf{p}$ yields a vector rotated with respect to $\mathbf{p}$ but of the same
norm, as follows directly from the orthogonality property
\eqref{eq:orthogonality}. A rotation matrix therefore admits three equivalent
geometrical meanings \cite{siciliano2009robotics}: it describes the mutual
orientation of two coordinate frames, it transforms the coordinates of a point
between two frames with a common origin, and it rotates a vector within one frame.

Rotations are combined by matrix multiplication. Considering three frames $0$, $1$
and $2$ with a common origin, applying \eqref{eq:coordinate_transform} twice gives
\begin{equation}
    ^0 R_2 = {}^0 R_1 \, ^1 R_2,
    \label{eq:rotation_composition}
\end{equation}
which is interpreted as a sequence of rotations, each defined with respect to the
frame resulting from the preceding one. The frame about whose axes a rotation is
made is termed the \textit{current frame}, and composition about current frames is
accordingly obtained by \textit{postmultiplication} in the order the rotations are
performed. If instead every rotation is referred to the initial frame, which is
then termed the \textit{fixed frame}, the same reasoning leads to composition by
\textit{premultiplication}. Since matrix multiplication is not commutative, two
rotations in general do not commute, and the resulting orientation depends on the
order in which they are applied. The Denavit-Hartenberg convention introduced
below is built on composition about current frames, and hence on postmultiplication
of the elementary rotations in \eqref{eq:elementary_rotations}.

\subsection{Homogeneous Transformation Matrix}

The rotation matrix and the translation vector can be combined into a single
description of \textit{pose} by means of the \textit{homogeneous transformation
matrix}, which allows the pose of an entire kinematic chain to be computed by
ordinary matrix multiplication. The pose of frame $B$
relative to frame $A$ is written as the $4 \times 4$ matrix
\begin{equation}
^A T_B = \begin{bmatrix}
n_x & s_x & a_x & P_x \\
n_y & s_y & a_y & P_y \\
n_z & s_z & a_z & P_z \\
0 & 0 & 0 & 1
\end{bmatrix} = \begin{bmatrix}
^A R_B & ^A P_B \\
\mathbf{0}^\top & 1
\end{bmatrix},
\end{equation}
where the upper-left $3 \times 3$ block $^A R_B$ is the rotation matrix describing
the relative orientation and the upper-right column $^A P_B$ is the position of the
origin of frame $B$ expressed in frame $A$. The columns of the rotation block are
frequently denoted $\mathbf{n}$, $\mathbf{s}$ and $\mathbf{a}$ (normal, sliding and
approach), a notation that originates from the description of gripper orientation.
The bottom row $\begin{bmatrix} 0 & 0 & 0 & 1 \end{bmatrix}$ has no geometric content
of its own but is what makes the composition of transformations reduce to matrix
multiplication.

The set of all such matrices forms the \textit{special Euclidean group} $SE(3)$, the
group of rigid-body motions in three-dimensional space. Its two central properties
are exactly the ones that make it the natural tool for serial manipulators. First,
transformations \textit{compose}: if the pose of frame $C$ is known relative to $B$,
and that of $B$ relative to $A$, then the pose of $C$ relative to $A$ is obtained
simply as
\begin{equation}
    ^A T_C = {}^A T_B \, ^B T_C,
\end{equation}
which generalises directly to a chain of arbitrarily many frames. Second, the inverse
transformation has the closed form
\begin{equation}
    ^A T_B^{-1} = {}^B T_A = \begin{bmatrix}
        ^A R_B^\top & -\,^A R_B^\top \, ^A P_B \\
        \mathbf{0}^\top & 1
    \end{bmatrix},
\end{equation}
so that changing the direction of a transformation requires no numerical matrix
inversion. The compositional property is the cornerstone of forward kinematics: by
attaching a frame to every link and writing the transformation between consecutive
frames, the pose of the end-effector relative to the base is recovered by multiplying
the individual link transformations together. The systematic assignment of these link
frames is governed by the Denavit-Hartenberg convention described next.

\subsection{The Denavit-Hartenberg (D-H) Convention}

The \acrshort{dh} convention is a standardised procedure for attaching coordinate
frames to the links of a serial manipulator. Without such a standard, the placement
of link frames would be arbitrary and the transformation between two consecutive
frames would in general require six parameters --- three for position and three for
orientation. The insight underlying the D-H convention is that, by choosing the frames
according to a fixed set of rules, the relative pose of two consecutive frames can
always be described by only \textit{four} parameters instead of six
\cite{siciliano2009robotics}. This reduction is achieved by aligning each frame's
$z$-axis with its joint axis and its $x$-axis along the common normal between
consecutive joint axes, which removes two of the degrees of freedom in the frame
placement.

The four parameters, illustrated for link $i$, are the \textit{link length} $a_i$,
the \textit{link twist} $\alpha_i$, the \textit{link offset} $d_i$, and the
\textit{joint angle} $\vartheta_i$. For a revolute joint the joint angle
$\vartheta_i$ is the variable quantity and the remaining three parameters are fixed
by the geometry of the link; for a prismatic joint the offset $d_i$ is variable and
the other three are constant.

The convention works along a recursive algorithm that transforms frame $i{-}1$ into
frame $i$ through four successive elementary transformations, performed in order:
\begin{enumerate}
    \item a rotation about the $z_{i-1}$-axis by the angle $\theta_i$, aligning
          $x_{i-1}$ with $x_i$;
    \item a translation along the $z_{i-1}$-axis by the offset $d_i$;
    \item a translation along the (now rotated) $x_i$-axis by the length $a_i$;
    \item a rotation about the $x_i$-axis by the twist $\alpha_i$, aligning
          $z_{i-1}$ with $z_i$.
\end{enumerate}
Expressing each step as a homogeneous transformation and multiplying them together
yields the link transformation matrix
\begin{equation}
^{i-1}T_i = \begin{bmatrix}
\cos\theta_i & -\sin\theta_i\cos\alpha_i & \sin\theta_i\sin\alpha_i & a_i\cos\theta_i \\
\sin\theta_i & \cos\theta_i\cos\alpha_i & -\cos\theta_i\sin\alpha_i & a_i\sin\theta_i \\
0 & \sin\alpha_i & \cos\alpha_i & d_i \\
0 & 0 & 0 & 1
\end{bmatrix}.
\label{eq:dh_matrix}
\end{equation}
Applying this procedure to every joint of the ViperX-300 yields the parameter set
listed in Table~\ref{tab:dh_params}, from which the complete kinematic model of the
manipulator follows.

\subsection{Forward and Inverse Kinematics}

\subsubsection{Forward Kinematics}
With a link frame assigned to each joint and the corresponding D-H transformation
\eqref{eq:dh_matrix} written for every link, the \textit{forward kinematics} of the
manipulator follows directly from the compositional property of homogeneous
transformations. Given a joint configuration
$\mathbf{q} = \begin{bmatrix} q_1 & \cdots & q_n \end{bmatrix}^\top$, the pose of the
end-effector frame $n$ relative to the base frame $0$ is obtained by chaining the
individual link transformations:
\begin{equation}
    ^0 T_n(\mathbf{q}) = {}^0 T_1(q_1)\, ^1 T_2(q_2) \;\cdots\; ^{n-1} T_n(q_n).
    \label{eq:fk_chain}
\end{equation}
The result is a single $4 \times 4$ matrix encoding both the position and the
orientation of the end-effector as a function of the joint variables. Forward
kinematics is, from a computational standpoint, entirely benign: for any given
$\mathbf{q}$ the pose is unique and is obtained by a fixed sequence of matrix
multiplications, with no iteration or ambiguity. It constitutes the mapping from
joint space to task space introduced earlier.

\subsubsection{Inverse Kinematics}
The \textit{inverse kinematics} problem reverses this mapping: given a desired
end-effector pose $^0 T_n^{\,\star}$, the task is to find a joint configuration
$\mathbf{q}$ for which $^0 T_n(\mathbf{q}) = {}^0 T_n^{\,\star}$. Although it is the
inverse of a problem that was straightforward, it is considerably more delicate, for
three reasons. First, \textit{existence} is not guaranteed: a solution exists only if
the requested pose lies within the reachable workspace of the manipulator. Second, the
solution is generally \textit{not unique}: a 6-DOF arm can frequently reach the same
pose through several distinct configurations --- the familiar elbow-up and elbow-down
postures being the classic example --- so the inverse map is one-to-many. Third, the
equations to be solved are \textit{nonlinear} in the joint variables, since the entries
of \eqref{eq:fk_chain} contain products of sines and cosines of the joint angles, and
for a general manipulator no closed-form solution exists.

In such cases the inverse kinematics must be solved numerically, typically by iterative
schemes based on the manipulator Jacobian, which raise the usual concerns of
convergence, computational cost and behaviour near singularities. Closed-form analytic
solutions are nonetheless available for manipulators whose geometry satisfies certain
conditions; the most important of these is the presence of a \textit{spherical wrist},
in which the axes of the final three joints intersect at a common point. This geometry
allows the problem to be \textit{decoupled}, with the wrist centre position determining
the first three joints and the end-effector orientation determining the last three,
each subproblem then admitting an analytic solution. Whether such a decoupling applies
to the ViperX-300 follows directly from its D-H parameters and is examined in the
kinematic analysis of this work.

% ============================================================

Robotics may come in many different forms but they all have something in
common: they share the same basics. The modern studies of robots include
everything from kinematics and trajectory generation to motion planning. In
order to grasp these more advanced concepts one must first recognise the
fundamentals. As a foundation of the subject, the following draws on
\cite{lynch2017modern}, which provides a clear explanation of the key concepts
of robotic manipulators.

\subsection{Configuration Space}
%When entering the mathematical way of explaining how manipulators move and  (Explaining joint space vs. task space)
To construct a robot you need various different components. The basic \textit{link} is rigid body, to permit movement of one link relative to the next one, a \textit{joint} is connected between the two consecutive links. At its core, a manipulator is a chain of such links coupled by joints. To realise the complex current pose of the manipulator, it can always be reduced to a couple of number describing the lengths of the links and the angle between them.
The complete specification of every position of each point of the robot is a \textit{configuration}. The minimum number of real-valued coordinates required to specify the configuration is the number of degrees of freedom of the robot, and the space of all attainable configurations is the \textit{configuration space}, of C-space \cite{lynch2017modern}. For a manipulator built from any number of $n$ joints adding one degree of freedom each, a configuration is conveniently collected intro the joint variable vector
\begin{equation}
    \mathbf{q} = \begin{bmatrix} q_1 & q_2 & \cdots & q_n \end{bmatrix}^\top \in \mathbb{R}^n,
\end{equation}
where each $q_i$ is a joint angle for a revolute joint or a linear displacement for prismatic joint.
When addressing the configuration space, generally the space is not configured to the conventional Euclidean space but rather a manifold. Taking a single revolute joint for example, it cant only move along a circle's outer rim and so forth its C-space is confined to a circle. The configuration space of a $n$-joint revolute arm will therefore result in an $n$-torus.

The configuration space described above, with each point given by the joint
variable vector $\mathbf{q}$, is what is referred to as the \textit{joint space}
of the manipulator. These are the variables that are directly driven by the
actuators and read back by the sensors. The joint space is however seldom the
space in which a task is actually described. When a manipulator is asked to reach
a point or to align its gripper, the task is expressed through the position and
orientation of the end-effector, which belong to a separate space called the
\textit{task space} (or operational space). Rather than joint values, the task
space is built from poses in ordinary three-dimensional space.

The link between the two spaces is the kinematics of the manipulator. The
\textit{forward kinematics} takes a set of joint values from the joint space and
returns the corresponding end-effector pose in the task space, while the
\textit{inverse kinematics} goes the other way, starting from a desired pose and
solving for the joint values that achieve it. Deriving and studying these two
mappings is the main focus of the sections that follow.

\subsection{Degrees of Freedom}

When specifying the \gls{dof} of a robot, one is specifying how the robot can
move in space. To start, the classic but straightforward example given in
\cite{lynch2017modern} explains the simple door fixed by hinges to the frame.
On its own, a rigid body moving freely in three-dimensional space has six
degrees of freedom. Attaching the door to the wall through the hinge introduces
constraints: the hinge permits rotation about its axis through an angle $\theta$
while removing every other possible motion. Five of the six degrees of freedom
are thereby constrained, leaving a single independent coordinate. The door
therefore has only one degree of freedom.

\begin{figure}[ht]
  \centering
  \input{figures/joints_dof.tex}
  \caption{The 6 different joint categories.}
  \label{fig:joints_dof}
\end{figure}

\subsection{Anatomy of an Open-Chain Robot}

The manipulators considered in this thesis are \textit{open-chain}, or serial,
mechanisms. In such a structure the links form a single unbranched sequence:
starting from a fixed \textit{base} that is rigidly attached to the environment,
each link is connected to the next by exactly one joint, terminating in the
\textit{end-effector}, the tool or gripper through which the robot interacts with
its surroundings. The defining property of an open chain is that there is one and
only one kinematic path from the base to the end-effector, in contrast to
closed-chain mechanisms where links form one or more loops. This single-path
structure is what makes the recursive construction of the kinematics in the
following sections possible.

The \textit{joints} are the elements that introduce motion into the chain, and
each one removes some of the relative freedom between the two links it connects
while permitting motion along the remaining directions. The two joint types
relevant here are the \textit{revolute} joint, which allows a pure rotation about
a fixed axis and is parameterised by an angle, and the \textit{prismatic} joint,
which allows a pure translation along a fixed axis and is parameterised by a
displacement. Both contribute a single degree of freedom. More general joints
such as cylindrical, spherical and universal joints exist and can be modelled as
combinations of these elementary joints; the six joint categories are illustrated
in Figure~\ref{fig:joints_dof}. The ViperX-300 studied in this work is a 6-DOF
open-chain manipulator constructed entirely from revolute joints, a configuration
that grants its end-effector six degrees of freedom and thus, in principle, the
ability to attain an arbitrary position \textit{and} orientation within its
reachable workspace.

Motion is produced by the \textit{actuators} that drive the joints. In modern
manipulators these are most commonly electric servomotors, typically brushless DC
or stepper motors combined with a gear reduction to trade speed for the torque
required to support the arm and its payload; hydraulic and pneumatic actuators are
also used where high force density or compliance is needed. Each actuator applies
a torque (for a revolute joint) or a force (for a prismatic joint) at its
respective joint in response to a commanded control signal. Because the links are
coupled, a torque applied at one joint propagates through the entire chain and
influences the motion of every link beyond it. Capturing this coupling
quantitatively is precisely the role of the dynamic model, which relates the
vector of joint torques to the resulting joint accelerations and is treated later
in this chapter.

\subsection{Rotation Matrix}

One of the fundamental tasks in modelling an open-chain robotic manipulator is to
describe the position and orientation of its end-effector. To do this, a
\textit{coordinate frame} is attached to each rigid body of the robot, giving a
systematic way of relating points and directions along the chain. The kinematics
then describes the geometric relations between these frames. To relate one frame
to another, two separate pieces of information are needed: the \textit{position}
of its origin, described by a translation vector, and its \textit{orientation},
described by a rotation matrix.

The orientation of a frame $B$ relative to a frame $A$ is conveniently described
by attaching an orthonormal frame to the body and expressing its unit vectors in
the coordinates of the reference frame \cite{siciliano2009robotics}. Collecting
the three unit vectors $\mathbf{x}_B$, $\mathbf{y}_B$ and $\mathbf{z}_B$ as
columns gives the $3 \times 3$ \textit{rotation matrix}
\begin{equation}
    ^A R_B = \begin{bmatrix}
        \mathbf{x}_B & \mathbf{y}_B & \mathbf{z}_B
    \end{bmatrix},
    \label{eq:rotation_matrix}
\end{equation}
whose entries are the direction cosines of the axes of frame $B$ with respect to
the axes of frame $A$.

Since the columns of $^A R_B$ are the unit vectors of an orthonormal frame, they
are mutually orthogonal and of unit norm. These conditions are compactly expressed
as
\begin{equation}
    ^A R_B^\top \, ^A R_B = I,
    \label{eq:orthogonality}
\end{equation}
that is, the rotation matrix is orthogonal. Postmultiplying \eqref{eq:orthogonality}
by the inverse matrix gives the useful result
\begin{equation}
    ^A R_B^{-1} = {}^A R_B^\top = {}^B R_A,
\end{equation}
so that reversing the roles of the two frames requires nothing more than a
transposition. The determinant of a rotation matrix equals $+1$ for a right-handed
frame and $-1$ for a left-handed one; since only right-handed frames are used here,
the matrices of interest belong to the \textit{special orthogonal group}
\begin{equation}
    SO(3) = \left\{ R \in \mathbb{R}^{3\times 3} \;\middle|\;
    R^\top R = I, \; \det R = +1 \right\}.
\end{equation}

Frames obtained by rotating the reference frame about one of its coordinate axes
are described by the \textit{elementary rotation matrices}. Taking the rotation
positive when made counter-clockwise about the axis in question, a rotation by an
angle $\vartheta$ about the $z$-, $y$- and $x$-axis is respectively given by
\begin{equation}
    R_z(\vartheta) = \begin{bmatrix}
        \cos\vartheta & -\sin\vartheta & 0 \\
        \sin\vartheta & \cos\vartheta & 0 \\
        0 & 0 & 1
    \end{bmatrix}, \quad
    R_y(\vartheta) = \begin{bmatrix}
        \cos\vartheta & 0 & \sin\vartheta \\
        0 & 1 & 0 \\
        -\sin\vartheta & 0 & \cos\vartheta
    \end{bmatrix}, \quad
    R_x(\vartheta) = \begin{bmatrix}
        1 & 0 & 0 \\
        0 & \cos\vartheta & -\sin\vartheta \\
        0 & \sin\vartheta & \cos\vartheta
    \end{bmatrix}.
    \label{eq:elementary_rotations}
\end{equation}
For these matrices $R_k(-\vartheta) = R_k^\top(\vartheta)$ holds for
$k = x, y, z$. They serve as the building blocks from which rotations about an
arbitrary axis in space are composed, and they are the ones that reappear in the
Denavit-Hartenberg convention.

Beyond describing orientation, the rotation matrix also relates the coordinates
of a point in two frames sharing a common origin. If a point $P$ is expressed as
$^B\mathbf{p}$ in frame $B$, its coordinates in frame $A$ follow from
\eqref{eq:rotation_matrix} as
\begin{equation}
    ^A\mathbf{p} = {}^A R_B \, ^B\mathbf{p},
    \label{eq:coordinate_transform}
\end{equation}
with the inverse transformation given by
$^B\mathbf{p} = {}^A R_B^\top \, ^A\mathbf{p}$. A third interpretation is obtained
by regarding $^A R_B$ as an operator acting within a single frame: the product
$R\mathbf{p}$ yields a vector rotated with respect to $\mathbf{p}$ but of the same
norm, as follows directly from the orthogonality property
\eqref{eq:orthogonality}. A rotation matrix therefore admits three equivalent
geometrical meanings \cite{siciliano2009robotics}: it describes the mutual
orientation of two coordinate frames, it transforms the coordinates of a point
between two frames with a common origin, and it rotates a vector within one frame.

Rotations are combined by matrix multiplication. Considering three frames $0$, $1$
and $2$ with a common origin, applying \eqref{eq:coordinate_transform} twice gives
\begin{equation}
    ^0 R_2 = {}^0 R_1 \, ^1 R_2,
    \label{eq:rotation_composition}
\end{equation}
which is interpreted as a sequence of rotations, each defined with respect to the
frame resulting from the preceding one. The frame about whose axes a rotation is
made is termed the \textit{current frame}, and composition about current frames is
accordingly obtained by \textit{postmultiplication} in the order the rotations are
performed. If instead every rotation is referred to the initial frame, which is
then termed the \textit{fixed frame}, the same reasoning leads to composition by
\textit{premultiplication}. Since matrix multiplication is not commutative, two
rotations in general do not commute, and the resulting orientation depends on the
order in which they are applied. The Denavit-Hartenberg convention introduced
below is built on composition about current frames, and hence on postmultiplication
of the elementary rotations in \eqref{eq:elementary_rotations}.

\subsection{Homogeneous Transformation Matrix}

The rotation matrix and the translation vector can be combined into a single
description of \textit{pose} by means of the \textit{homogeneous transformation
matrix}, which allows the pose of an entire kinematic chain to be computed by
ordinary matrix multiplication. The pose of frame $B$
relative to frame $A$ is written as the $4 \times 4$ matrix
\begin{equation}
^A T_B = \begin{bmatrix}
n_x & s_x & a_x & P_x \\
n_y & s_y & a_y & P_y \\
n_z & s_z & a_z & P_z \\
0 & 0 & 0 & 1
\end{bmatrix} = \begin{bmatrix}
^A R_B & ^A P_B \\
\mathbf{0}^\top & 1
\end{bmatrix},
\end{equation}
where the upper-left $3 \times 3$ block $^A R_B$ is the rotation matrix describing
the relative orientation and the upper-right column $^A P_B$ is the position of the
origin of frame $B$ expressed in frame $A$. The columns of the rotation block are
frequently denoted $\mathbf{n}$, $\mathbf{s}$ and $\mathbf{a}$ (normal, sliding and
approach), a notation that originates from the description of gripper orientation.
The bottom row $\begin{bmatrix} 0 & 0 & 0 & 1 \end{bmatrix}$ has no geometric content
of its own but is what makes the composition of transformations reduce to matrix
multiplication.

The set of all such matrices forms the \textit{special Euclidean group} $SE(3)$, the
group of rigid-body motions in three-dimensional space. Its two central properties
are exactly the ones that make it the natural tool for serial manipulators. First,
transformations \textit{compose}: if the pose of frame $C$ is known relative to $B$,
and that of $B$ relative to $A$, then the pose of $C$ relative to $A$ is obtained
simply as
\begin{equation}
    ^A T_C = {}^A T_B \, ^B T_C,
\end{equation}
which generalises directly to a chain of arbitrarily many frames. Second, the inverse
transformation has the closed form
\begin{equation}
    ^A T_B^{-1} = {}^B T_A = \begin{bmatrix}
        ^A R_B^\top & -\,^A R_B^\top \, ^A P_B \\
        \mathbf{0}^\top & 1
    \end{bmatrix},
\end{equation}
so that changing the direction of a transformation requires no numerical matrix
inversion. The compositional property is the cornerstone of forward kinematics: by
attaching a frame to every link and writing the transformation between consecutive
frames, the pose of the end-effector relative to the base is recovered by multiplying
the individual link transformations together. The systematic assignment of these link
frames is governed by the Denavit-Hartenberg convention described next.

\subsection{The Denavit-Hartenberg (D-H) Convention}

The \acrshort{dh} convention is a standardised procedure for attaching coordinate
frames to the links of a serial manipulator. Without such a standard, the placement
of link frames would be arbitrary and the transformation between two consecutive
frames would in general require six parameters --- three for position and three for
orientation. The insight underlying the D-H convention is that, by choosing the frames
according to a fixed set of rules, the relative pose of two consecutive frames can
always be described by only \textit{four} parameters instead of six
\cite{siciliano2009robotics}. This reduction is achieved by aligning each frame's
$z$-axis with its joint axis and its $x$-axis along the common normal between
consecutive joint axes, which removes two of the degrees of freedom in the frame
placement.

The four parameters, illustrated for link $i$, are the \textit{link length} $a_i$,
the \textit{link twist} $\alpha_i$, the \textit{link offset} $d_i$, and the
\textit{joint angle} $\vartheta_i$. For a revolute joint the joint angle
$\vartheta_i$ is the variable quantity and the remaining three parameters are fixed
by the geometry of the link; for a prismatic joint the offset $d_i$ is variable and
the other three are constant.

The convention works along a recursive algorithm that transforms frame $i{-}1$ into
frame $i$ through four successive elementary transformations, performed in order:
\begin{enumerate}
    \item a rotation about the $z_{i-1}$-axis by the angle $\theta_i$, aligning
          $x_{i-1}$ with $x_i$;
    \item a translation along the $z_{i-1}$-axis by the offset $d_i$;
    \item a translation along the (now rotated) $x_i$-axis by the length $a_i$;
    \item a rotation about the $x_i$-axis by the twist $\alpha_i$, aligning
          $z_{i-1}$ with $z_i$.
\end{enumerate}
Expressing each step as a homogeneous transformation and multiplying them together
yields the link transformation matrix
\begin{equation}
^{i-1}T_i = \begin{bmatrix}
\cos\theta_i & -\sin\theta_i\cos\alpha_i & \sin\theta_i\sin\alpha_i & a_i\cos\theta_i \\
\sin\theta_i & \cos\theta_i\cos\alpha_i & -\cos\theta_i\sin\alpha_i & a_i\sin\theta_i \\
0 & \sin\alpha_i & \cos\alpha_i & d_i \\
0 & 0 & 0 & 1
\end{bmatrix}.
\label{eq:dh_matrix}
\end{equation}
Applying this procedure to every joint of the ViperX-300 yields the parameter set
listed in Table~\ref{tab:dh_params}, from which the complete kinematic model of the
manipulator follows.

\subsection{Forward and Inverse Kinematics}

\subsubsection{Forward Kinematics}
With a link frame assigned to each joint and the corresponding D-H transformation
\eqref{eq:dh_matrix} written for every link, the \textit{forward kinematics} of the
manipulator follows directly from the compositional property of homogeneous
transformations. Given a joint configuration
$\mathbf{q} = \begin{bmatrix} q_1 & \cdots & q_n \end{bmatrix}^\top$, the pose of the
end-effector frame $n$ relative to the base frame $0$ is obtained by chaining the
individual link transformations:
\begin{equation}
    ^0 T_n(\mathbf{q}) = {}^0 T_1(q_1)\, ^1 T_2(q_2) \;\cdots\; ^{n-1} T_n(q_n).
    \label{eq:fk_chain}
\end{equation}
The result is a single $4 \times 4$ matrix encoding both the position and the
orientation of the end-effector as a function of the joint variables. Forward
kinematics is, from a computational standpoint, entirely benign: for any given
$\mathbf{q}$ the pose is unique and is obtained by a fixed sequence of matrix
multiplications, with no iteration or ambiguity. It constitutes the mapping from
joint space to task space introduced earlier.

\subsubsection{Inverse Kinematics}
The \textit{inverse kinematics} problem reverses this mapping: given a desired
end-effector pose $^0 T_n^{\,\star}$, the task is to find a joint configuration
$\mathbf{q}$ for which $^0 T_n(\mathbf{q}) = {}^0 T_n^{\,\star}$. Although it is the
inverse of a problem that was straightforward, it is considerably more delicate, for
three reasons. First, \textit{existence} is not guaranteed: a solution exists only if
the requested pose lies within the reachable workspace of the manipulator. Second, the
solution is generally \textit{not unique}: a 6-DOF arm can frequently reach the same
pose through several distinct configurations --- the familiar elbow-up and elbow-down
postures being the classic example --- so the inverse map is one-to-many. Third, the
equations to be solved are \textit{nonlinear} in the joint variables, since the entries
of \eqref{eq:fk_chain} contain products of sines and cosines of the joint angles, and
for a general manipulator no closed-form solution exists.

In such cases the inverse kinematics must be solved numerically, typically by iterative
schemes based on the manipulator Jacobian, which raise the usual concerns of
convergence, computational cost and behaviour near singularities. Closed-form analytic
solutions are nonetheless available for manipulators whose geometry satisfies certain
conditions; the most important of these is the presence of a \textit{spherical wrist},
in which the axes of the final three joints intersect at a common point. This geometry
allows the problem to be \textit{decoupled}, with the wrist centre position determining
the first three joints and the end-effector orientation determining the last three,
each subproblem then admitting an analytic solution. Whether such a decoupling applies
to the ViperX-300 follows directly from its D-H parameters and is examined in the
kinematic analysis of this work.

% ============================================================

Robotics may come in many different forms but they all have something in
common: they share the same basics. The modern studies of robots include
everything from kinematics and trajectory generation to motion planning. In
order to grasp these more advanced concepts one must first recognise the
fundamentals. As a foundation of the subject, the following draws on
\cite{lynch2017modern}, which provides a clear explanation of the key concepts
of robotic manipulators.

\subsection{Configuration Space}
%When entering the mathematical way of explaining how manipulators move and  (Explaining joint space vs. task space)
To construct a robot you need various different components. The basic \textit{link} is rigid body, to permit movement of one link relative to the next one, a \textit{joint} is connected between the two consecutive links. At its core, a manipulator is a chain of such links coupled by joints. To realise the complex current pose of the manipulator, it can always be reduced to a couple of number describing the lengths of the links and the angle between them.
The complete specification of every position of each point of the robot is a \textit{configuration}. The minimum number of real-valued coordinates required to specify the configuration is the number of degrees of freedom of the robot, and the space of all attainable configurations is the \textit{configuration space}, of C-space \cite{lynch2017modern}. For a manipulator built from any number of $n$ joints adding one degree of freedom each, a configuration is conveniently collected intro the joint variable vector
\begin{equation}
    \mathbf{q} = \begin{bmatrix} q_1 & q_2 & \cdots & q_n \end{bmatrix}^\top \in \mathbb{R}^n,
\end{equation}
where each $q_i$ is a joint angle for a revolute joint or a linear displacement for prismatic joint.
When addressing the configuration space, generally the space is not configured to the conventional Euclidean space but rather a manifold. Taking a single revolute joint for example, it cant only move along a circle's outer rim and so forth its C-space is confined to a circle. The configuration space of a $n$-joint revolute arm will therefore result in an $n$-torus.

The configuration space described above, with each point given by the joint
variable vector $\mathbf{q}$, is what is referred to as the \textit{joint space}
of the manipulator. These are the variables that are directly driven by the
actuators and read back by the sensors. The joint space is however seldom the
space in which a task is actually described. When a manipulator is asked to reach
a point or to align its gripper, the task is expressed through the position and
orientation of the end-effector, which belong to a separate space called the
\textit{task space} (or operational space). Rather than joint values, the task
space is built from poses in ordinary three-dimensional space.

The link between the two spaces is the kinematics of the manipulator. The
\textit{forward kinematics} takes a set of joint values from the joint space and
returns the corresponding end-effector pose in the task space, while the
\textit{inverse kinematics} goes the other way, starting from a desired pose and
solving for the joint values that achieve it. Deriving and studying these two
mappings is the main focus of the sections that follow.

\subsection{Degrees of Freedom}

When specifying the \gls{dof} of a robot, one is specifying how the robot can
move in space. To start, the classic but straightforward example given in
\cite{lynch2017modern} explains the simple door fixed by hinges to the frame.
On its own, a rigid body moving freely in three-dimensional space has six
degrees of freedom. Attaching the door to the wall through the hinge introduces
constraints: the hinge permits rotation about its axis through an angle $\theta$
while removing every other possible motion. Five of the six degrees of freedom
are thereby constrained, leaving a single independent coordinate. The door
therefore has only one degree of freedom.

\begin{figure}[ht]
  \centering
  \input{figures/joints_dof.tex}
  \caption{The 6 different joint categories.}
  \label{fig:joints_dof}
\end{figure}

\subsection{Anatomy of an Open-Chain Robot}

The manipulators considered in this thesis are \textit{open-chain}, or serial,
mechanisms. In such a structure the links form a single unbranched sequence:
starting from a fixed \textit{base} that is rigidly attached to the environment,
each link is connected to the next by exactly one joint, terminating in the
\textit{end-effector}, the tool or gripper through which the robot interacts with
its surroundings. The defining property of an open chain is that there is one and
only one kinematic path from the base to the end-effector, in contrast to
closed-chain mechanisms where links form one or more loops. This single-path
structure is what makes the recursive construction of the kinematics in the
following sections possible.

The \textit{joints} are the elements that introduce motion into the chain, and
each one removes some of the relative freedom between the two links it connects
while permitting motion along the remaining directions. The two joint types
relevant here are the \textit{revolute} joint, which allows a pure rotation about
a fixed axis and is parameterised by an angle, and the \textit{prismatic} joint,
which allows a pure translation along a fixed axis and is parameterised by a
displacement. Both contribute a single degree of freedom. More general joints
such as cylindrical, spherical and universal joints exist and can be modelled as
combinations of these elementary joints; the six joint categories are illustrated
in Figure~\ref{fig:joints_dof}. The ViperX-300 studied in this work is a 6-DOF
open-chain manipulator constructed entirely from revolute joints, a configuration
that grants its end-effector six degrees of freedom and thus, in principle, the
ability to attain an arbitrary position \textit{and} orientation within its
reachable workspace.

Motion is produced by the \textit{actuators} that drive the joints. In modern
manipulators these are most commonly electric servomotors, typically brushless DC
or stepper motors combined with a gear reduction to trade speed for the torque
required to support the arm and its payload; hydraulic and pneumatic actuators are
also used where high force density or compliance is needed. Each actuator applies
a torque (for a revolute joint) or a force (for a prismatic joint) at its
respective joint in response to a commanded control signal. Because the links are
coupled, a torque applied at one joint propagates through the entire chain and
influences the motion of every link beyond it. Capturing this coupling
quantitatively is precisely the role of the dynamic model, which relates the
vector of joint torques to the resulting joint accelerations and is treated later
in this chapter.

\subsection{Rotation Matrix}

One of the fundamental tasks in modelling an open-chain robotic manipulator is to
describe the position and orientation of its end-effector. To do this, a
\textit{coordinate frame} is attached to each rigid body of the robot, giving a
systematic way of relating points and directions along the chain. The kinematics
then describes the geometric relations between these frames. To relate one frame
to another, two separate pieces of information are needed: the \textit{position}
of its origin, described by a translation vector, and its \textit{orientation},
described by a rotation matrix.

The orientation of a frame $B$ relative to a frame $A$ is conveniently described
by attaching an orthonormal frame to the body and expressing its unit vectors in
the coordinates of the reference frame \cite{siciliano2009robotics}. Collecting
the three unit vectors $\mathbf{x}_B$, $\mathbf{y}_B$ and $\mathbf{z}_B$ as
columns gives the $3 \times 3$ \textit{rotation matrix}
\begin{equation}
    ^A R_B = \begin{bmatrix}
        \mathbf{x}_B & \mathbf{y}_B & \mathbf{z}_B
    \end{bmatrix},
    \label{eq:rotation_matrix}
\end{equation}
whose entries are the direction cosines of the axes of frame $B$ with respect to
the axes of frame $A$.

Since the columns of $^A R_B$ are the unit vectors of an orthonormal frame, they
are mutually orthogonal and of unit norm. These conditions are compactly expressed
as
\begin{equation}
    ^A R_B^\top \, ^A R_B = I,
    \label{eq:orthogonality}
\end{equation}
that is, the rotation matrix is orthogonal. Postmultiplying \eqref{eq:orthogonality}
by the inverse matrix gives the useful result
\begin{equation}
    ^A R_B^{-1} = {}^A R_B^\top = {}^B R_A,
\end{equation}
so that reversing the roles of the two frames requires nothing more than a
transposition. The determinant of a rotation matrix equals $+1$ for a right-handed
frame and $-1$ for a left-handed one; since only right-handed frames are used here,
the matrices of interest belong to the \textit{special orthogonal group}
\begin{equation}
    SO(3) = \left\{ R \in \mathbb{R}^{3\times 3} \;\middle|\;
    R^\top R = I, \; \det R = +1 \right\}.
\end{equation}

Frames obtained by rotating the reference frame about one of its coordinate axes
are described by the \textit{elementary rotation matrices}. Taking the rotation
positive when made counter-clockwise about the axis in question, a rotation by an
angle $\vartheta$ about the $z$-, $y$- and $x$-axis is respectively given by
\begin{equation}
    R_z(\vartheta) = \begin{bmatrix}
        \cos\vartheta & -\sin\vartheta & 0 \\
        \sin\vartheta & \cos\vartheta & 0 \\
        0 & 0 & 1
    \end{bmatrix}, \quad
    R_y(\vartheta) = \begin{bmatrix}
        \cos\vartheta & 0 & \sin\vartheta \\
        0 & 1 & 0 \\
        -\sin\vartheta & 0 & \cos\vartheta
    \end{bmatrix}, \quad
    R_x(\vartheta) = \begin{bmatrix}
        1 & 0 & 0 \\
        0 & \cos\vartheta & -\sin\vartheta \\
        0 & \sin\vartheta & \cos\vartheta
    \end{bmatrix}.
    \label{eq:elementary_rotations}
\end{equation}
For these matrices $R_k(-\vartheta) = R_k^\top(\vartheta)$ holds for
$k = x, y, z$. They serve as the building blocks from which rotations about an
arbitrary axis in space are composed, and they are the ones that reappear in the
Denavit-Hartenberg convention.

Beyond describing orientation, the rotation matrix also relates the coordinates
of a point in two frames sharing a common origin. If a point $P$ is expressed as
$^B\mathbf{p}$ in frame $B$, its coordinates in frame $A$ follow from
\eqref{eq:rotation_matrix} as
\begin{equation}
    ^A\mathbf{p} = {}^A R_B \, ^B\mathbf{p},
    \label{eq:coordinate_transform}
\end{equation}
with the inverse transformation given by
$^B\mathbf{p} = {}^A R_B^\top \, ^A\mathbf{p}$. A third interpretation is obtained
by regarding $^A R_B$ as an operator acting within a single frame: the product
$R\mathbf{p}$ yields a vector rotated with respect to $\mathbf{p}$ but of the same
norm, as follows directly from the orthogonality property
\eqref{eq:orthogonality}. A rotation matrix therefore admits three equivalent
geometrical meanings \cite{siciliano2009robotics}: it describes the mutual
orientation of two coordinate frames, it transforms the coordinates of a point
between two frames with a common origin, and it rotates a vector within one frame.

Rotations are combined by matrix multiplication. Considering three frames $0$, $1$
and $2$ with a common origin, applying \eqref{eq:coordinate_transform} twice gives
\begin{equation}
    ^0 R_2 = {}^0 R_1 \, ^1 R_2,
    \label{eq:rotation_composition}
\end{equation}
which is interpreted as a sequence of rotations, each defined with respect to the
frame resulting from the preceding one. The frame about whose axes a rotation is
made is termed the \textit{current frame}, and composition about current frames is
accordingly obtained by \textit{postmultiplication} in the order the rotations are
performed. If instead every rotation is referred to the initial frame, which is
then termed the \textit{fixed frame}, the same reasoning leads to composition by
\textit{premultiplication}. Since matrix multiplication is not commutative, two
rotations in general do not commute, and the resulting orientation depends on the
order in which they are applied. The Denavit-Hartenberg convention introduced
below is built on composition about current frames, and hence on postmultiplication
of the elementary rotations in \eqref{eq:elementary_rotations}.

\subsection{Homogeneous Transformation Matrix}

The rotation matrix and the translation vector can be combined into a single
description of \textit{pose} by means of the \textit{homogeneous transformation
matrix}, which allows the pose of an entire kinematic chain to be computed by
ordinary matrix multiplication. The pose of frame $B$
relative to frame $A$ is written as the $4 \times 4$ matrix
\begin{equation}
^A T_B = \begin{bmatrix}
n_x & s_x & a_x & P_x \\
n_y & s_y & a_y & P_y \\
n_z & s_z & a_z & P_z \\
0 & 0 & 0 & 1
\end{bmatrix} = \begin{bmatrix}
^A R_B & ^A P_B \\
\mathbf{0}^\top & 1
\end{bmatrix},
\end{equation}
where the upper-left $3 \times 3$ block $^A R_B$ is the rotation matrix describing
the relative orientation and the upper-right column $^A P_B$ is the position of the
origin of frame $B$ expressed in frame $A$. The columns of the rotation block are
frequently denoted $\mathbf{n}$, $\mathbf{s}$ and $\mathbf{a}$ (normal, sliding and
approach), a notation that originates from the description of gripper orientation.
The bottom row $\begin{bmatrix} 0 & 0 & 0 & 1 \end{bmatrix}$ has no geometric content
of its own but is what makes the composition of transformations reduce to matrix
multiplication.

The set of all such matrices forms the \textit{special Euclidean group} $SE(3)$, the
group of rigid-body motions in three-dimensional space. Its two central properties
are exactly the ones that make it the natural tool for serial manipulators. First,
transformations \textit{compose}: if the pose of frame $C$ is known relative to $B$,
and that of $B$ relative to $A$, then the pose of $C$ relative to $A$ is obtained
simply as
\begin{equation}
    ^A T_C = {}^A T_B \, ^B T_C,
\end{equation}
which generalises directly to a chain of arbitrarily many frames. Second, the inverse
transformation has the closed form
\begin{equation}
    ^A T_B^{-1} = {}^B T_A = \begin{bmatrix}
        ^A R_B^\top & -\,^A R_B^\top \, ^A P_B \\
        \mathbf{0}^\top & 1
    \end{bmatrix},
\end{equation}
so that changing the direction of a transformation requires no numerical matrix
inversion. The compositional property is the cornerstone of forward kinematics: by
attaching a frame to every link and writing the transformation between consecutive
frames, the pose of the end-effector relative to the base is recovered by multiplying
the individual link transformations together. The systematic assignment of these link
frames is governed by the Denavit-Hartenberg convention described next.

\subsection{The Denavit-Hartenberg (D-H) Convention}

The \acrshort{dh} convention is a standardised procedure for attaching coordinate
frames to the links of a serial manipulator. Without such a standard, the placement
of link frames would be arbitrary and the transformation between two consecutive
frames would in general require six parameters --- three for position and three for
orientation. The insight underlying the D-H convention is that, by choosing the frames
according to a fixed set of rules, the relative pose of two consecutive frames can
always be described by only \textit{four} parameters instead of six
\cite{siciliano2009robotics}. This reduction is achieved by aligning each frame's
$z$-axis with its joint axis and its $x$-axis along the common normal between
consecutive joint axes, which removes two of the degrees of freedom in the frame
placement.

The four parameters, illustrated for link $i$, are the \textit{link length} $a_i$,
the \textit{link twist} $\alpha_i$, the \textit{link offset} $d_i$, and the
\textit{joint angle} $\vartheta_i$. For a revolute joint the joint angle
$\vartheta_i$ is the variable quantity and the remaining three parameters are fixed
by the geometry of the link; for a prismatic joint the offset $d_i$ is variable and
the other three are constant.

The convention works along a recursive algorithm that transforms frame $i{-}1$ into
frame $i$ through four successive elementary transformations, performed in order:
\begin{enumerate}
    \item a rotation about the $z_{i-1}$-axis by the angle $\theta_i$, aligning
          $x_{i-1}$ with $x_i$;
    \item a translation along the $z_{i-1}$-axis by the offset $d_i$;
    \item a translation along the (now rotated) $x_i$-axis by the length $a_i$;
    \item a rotation about the $x_i$-axis by the twist $\alpha_i$, aligning
          $z_{i-1}$ with $z_i$.
\end{enumerate}
Expressing each step as a homogeneous transformation and multiplying them together
yields the link transformation matrix
\begin{equation}
^{i-1}T_i = \begin{bmatrix}
\cos\theta_i & -\sin\theta_i\cos\alpha_i & \sin\theta_i\sin\alpha_i & a_i\cos\theta_i \\
\sin\theta_i & \cos\theta_i\cos\alpha_i & -\cos\theta_i\sin\alpha_i & a_i\sin\theta_i \\
0 & \sin\alpha_i & \cos\alpha_i & d_i \\
0 & 0 & 0 & 1
\end{bmatrix}.
\label{eq:dh_matrix}
\end{equation}
Applying this procedure to every joint of the ViperX-300 yields the parameter set
listed in Table~\ref{tab:dh_params}, from which the complete kinematic model of the
manipulator follows.

\subsection{Forward and Inverse Kinematics}

\subsubsection{Forward Kinematics}
With a link frame assigned to each joint and the corresponding D-H transformation
\eqref{eq:dh_matrix} written for every link, the \textit{forward kinematics} of the
manipulator follows directly from the compositional property of homogeneous
transformations. Given a joint configuration
$\mathbf{q} = \begin{bmatrix} q_1 & \cdots & q_n \end{bmatrix}^\top$, the pose of the
end-effector frame $n$ relative to the base frame $0$ is obtained by chaining the
individual link transformations:
\begin{equation}
    ^0 T_n(\mathbf{q}) = {}^0 T_1(q_1)\, ^1 T_2(q_2) \;\cdots\; ^{n-1} T_n(q_n).
    \label{eq:fk_chain}
\end{equation}
The result is a single $4 \times 4$ matrix encoding both the position and the
orientation of the end-effector as a function of the joint variables. Forward
kinematics is, from a computational standpoint, entirely benign: for any given
$\mathbf{q}$ the pose is unique and is obtained by a fixed sequence of matrix
multiplications, with no iteration or ambiguity. It constitutes the mapping from
joint space to task space introduced earlier.

\subsubsection{Inverse Kinematics}
The \textit{inverse kinematics} problem reverses this mapping: given a desired
end-effector pose $^0 T_n^{\,\star}$, the task is to find a joint configuration
$\mathbf{q}$ for which $^0 T_n(\mathbf{q}) = {}^0 T_n^{\,\star}$. Although it is the
inverse of a problem that was straightforward, it is considerably more delicate, for
three reasons. First, \textit{existence} is not guaranteed: a solution exists only if
the requested pose lies within the reachable workspace of the manipulator. Second, the
solution is generally \textit{not unique}: a 6-DOF arm can frequently reach the same
pose through several distinct configurations --- the familiar elbow-up and elbow-down
postures being the classic example --- so the inverse map is one-to-many. Third, the
equations to be solved are \textit{nonlinear} in the joint variables, since the entries
of \eqref{eq:fk_chain} contain products of sines and cosines of the joint angles, and
for a general manipulator no closed-form solution exists.

In such cases the inverse kinematics must be solved numerically, typically by iterative
schemes based on the manipulator Jacobian, which raise the usual concerns of
convergence, computational cost and behaviour near singularities. Closed-form analytic
solutions are nonetheless available for manipulators whose geometry satisfies certain
conditions; the most important of these is the presence of a \textit{spherical wrist},
in which the axes of the final three joints intersect at a common point. This geometry
allows the problem to be \textit{decoupled}, with the wrist centre position determining
the first three joints and the end-effector orientation determining the last three,
each subproblem then admitting an analytic solution. Whether such a decoupling applies
to the ViperX-300 follows directly from its D-H parameters and is examined in the
kinematic analysis of this work.
